\documentclass[11pt]{article}
\usepackage[margin=25mm]{geometry}
\usepackage[T1]{fontenc}
\usepackage{amsmath,amssymb,booktabs,array}
\usepackage{hyperref}
\title{Native Material Profiles and Qualified Graph Readouts}
\author{}
\date{}
\begin{document}
\maketitle

\paragraph{Claims and presentations.}
Unqualified \texttt{set:} claims use the proved \texttt{prime-common}
material fragment. Additional logical, foundation, choice and universe laws
require an explicit profile or declaration. A graph presentation is data,
not an adoption of its interpretation's complete set theory. Profiles and
readouts are separate: the same stored presentation can be inspected under
several qualified readouts without changing its authored data.

The profile-indexed first-order source signature contains equality and
membership. A derivation retains its declarations, hypothesis occurrences
and law origins. Checked interpretations connect these derivations to the
constructed contextual graph model and to separately qualified material
models. The profile catalogue keeps four statuses separate: declared theory,
constructed model, checked interpretation and executable fragment. A
literature result, a finite picture or a native replay is not promoted to a
proof of a full native set theory.
Catalogue entries with different signatures are separate theories. A
comparison with the common material fragment requires a proved interpretation;
catalogue membership does not make an entry a stronger extension of that
fragment.

\paragraph{Finite values.}
A source \texttt{(graph r (row$_0$\ldots row$_{n-1}$))} has a valid root and
valid successor indices. Rows retain order and repeated occurrence positions.
The adjacency relation uses the support of those rows: a repeated index is
one edge, while two distinct indices remain two target vertices until the
selected material interpretation identifies them.

\begin{center}
\begin{tabular}{@{}>{\raggedright\arraybackslash}p{31mm}
                  >{\raggedright\arraybackslash}p{72mm}
                  >{\raggedright\arraybackslash}p{43mm}@{}}
\toprule
Readout & Identity identified & Boundary \\
\midrule
\texttt{presentation} & Authored graph data, including row occurrences & Structural equality \\
\texttt{aczel} & Greatest bisimulation; equality of the Aczel decorations & Material membership preserved \\
\texttt{raw-unfolding} & Isomorphism of full rooted unfolding trees & Distinct adjacency targets are counted \\
\texttt{scott-canonical} & Equality of projected nodes after iterated unfolding quotient and edge deduplication & Final graph is Scott-extensional \\
\texttt{foundation} & Aczel identity restricted to accessible roots & Partial, membership-preserving reading into well-founded sets \\
\texttt{observed} & Unlabelled graph behavior, outcome, scope and origin ledger & Results and faults remain distinct \\
\bottomrule
\end{tabular}
\end{center}

The canonical finite Scott construction computes actual quotient carriers,
surjective projections and member images. Every genuine merge decreases the
carrier size; the input size bounds the number of quotient rounds. Raw
unfolding comparison must be repeated after removing duplicate quotient
edges. For example, two empty child vertices and one empty child vertex
initially have different raw trees but the same canonical material reading.
Transport to another carrier requires a canonical decoration of the final
graph. Ordinary decorations alone do not determine this identity kernel.
Closed embeddings preserve complete successor fibres and induce closed
embeddings through the quotient iterations. Consequently a comparison in
a disjoint joint graph agrees with full unfolding isomorphism of the
independently canonical roots. Equality is reflexive, symmetric and
transitive, and adding unrelated graph components does not change it.
Membership agrees with equality to one of the parent's selected child
representatives. An injective edge-preserving map that introduces new
successors does not satisfy this preservation contract.
This finite construction does not supply a complete Scott anti-foundation
universe.

The Foundation reading checks accessibility of the root, rather than absence
of every cycle in the stored graph. An unreachable cycle is harmless to that
root; a reachable cycle is outside the readout's domain. No total operation
silently sends a cyclic set to the empty set and then claims to preserve
membership.

\paragraph{Evidence consumers.}
Material support may identify different presentations or different child
occurrences. A consumer using membership support can descend through that
quotient when its family and term are compatible with the proved kernel and
context transport. A consumer inspecting origin positions or provenance may
require the retained presentation instead. Compatible dependent sections
and their substitution comparisons are proved in the contextual model;
bare equality of displayed values does not establish those contracts.

An observed source is
\texttt{(observed-graph G outcome scope origins)}, where the outcome is
\texttt{(returned k)} or \texttt{(faulted k)}. Its origin receipt retains
the complete authored source, an in-range index and the corresponding
origin. Equality of observed readings permits transport of the appropriate
index and origin value; it does not replace the retained source by another
presentation. A proof-path count is a separate observation and is not a
count of independent evidence.
The finite observed surface does not itself encode labelled execution
histories or every future operational observation. Those require their
declared source semantics and the corresponding richer interpretation.

\paragraph{Proof compilation and authority.}
The general source compiler constructs retained typed HOL proof syntax.
Its codec and decoder check object-variable indices, hypothesis positions,
types and conclusions. Recovery is exact for the retained HOL target proof
tree, including its annotations and hypothesis occurrences. The first-order
compiler interprets weakening through an occurrence map and expands derived
connectives; its injection on the original source trees is not asserted.
The typed receipt envelope separately retains declarations, adoptions,
origins, source context and variable interpretation. Bytes alone do not
recover opaque primitive adoptions or authorize them. The native code
renderer is a separate boundary. Structural logical equality uses the
explicit \texttt{hol-equality} rule bundle; transport-only
\texttt{eq-subst} does not imply reflexivity. Native equality proofs lower
to dependent reflexivity and transport terms, which the regular kernel
checks in their complete admitted context. Neither a type annotation nor
serialized data adopts an assumption. Withdrawing a required declaration
invalidates reuse of a dependent theorem or saved native package.

\paragraph{Definitions and conservative extension.}
A finite history adds each fresh symbol with a closed body checked in the
preceding signature. The general dependent calculus expands the bodies in
actual typing, equality and subsumption derivations, including dependent
contexts, substitution and reindexing. Old sequents and old-type inhabitation
are conservative. Normalization inheritance additionally requires the
corresponding normalization theorem for the base calculus.

For fixed-simple-type HOL, each stage constructs a Henkin model from the
preceding model and the defining body. Its open-term denotation agrees with
independent syntactic expansion; model composition agrees with history
composition. Expansion traverses all 29 retained HOL rules, preserving the
ordered rule tree and each local hypothesis position. Every expanded
definition equation has a constructed reflexivity proof. Substituting these
proofs discharges the added equations, so an old sequent is derivable in the
extended calculus exactly when it was derivable before. Discharging equations
uses constructed reflexivity leaves. Substitution of arbitrary hypothesis
proofs can add branches. Neither a proof size nor its preservation implies
independence of the evidence.

Native function admission exposes only the successive function-type heads of
a transparent type alias; it leaves computations in dependent results for
their own checks. The authored type and its current dependency receipts remain
retained. A computationally equivalent replacement of a required definition
still changes its authority. These construction theorems and qualified native
controls do not establish general C-checker correctness or admission of
arbitrary recursive definitions.

\paragraph{Executable surface and qualification.}
\begin{verbatim}
(set:interpret finite-graph aczel (equal G H))
(set:interpret finite-graph aczel (member G H))
(set:interpret finite-graph scott-canonical (members G))
(set:interpret finite-graph foundation (eligible G))
(set:interpret finite-graph observed (origin-at O i))
\end{verbatim}
Equality and membership always name a readout. Returned finite values retain
their original presentations; ordinary atom equality is not material
equality. A list of material members returns first authored representatives,
not a canonical printing order. Known malformed closed input is refuted;
open input or a root outside the Foundation domain is undetermined.
Exhausted work or storage is incomplete. Unknown interpreters remain data.

General solver, interpretation, substitution and recovery theorems are
kernel evidence. Native execution is qualified separately against checked
solver output, source-proof exports and independent controls, with source,
artifact and binary hashes retained. This establishes the stated replay
agreement, not a proof of the C implementation's correctness for all inputs.
Native universal construction, larger model comparisons and unrestricted
library transport retain their own obligations.
\end{document}
